% TOP_PROBLEM: {"problemId": "problem.hilbert-s-third-problem-for-spherical-and-hyperbolic-scissors-congruence", "canonicalTitle": "Hilbert's Third Problem for Spherical and Hyperbolic Scissors Congruence", "releaseRank": 369, "primaryDomain": "geometry_topology", "primaryDomainLabel": "Geometry & topology", "displayStatus": "open", "releaseStatus": "open"}
% !TeX program = lualatex
% ATTEMPT_STATUS: unresolved
\documentclass[11pt]{article}
\usepackage[margin=25mm]{geometry}
\usepackage{fontspec}
\setmainfont{Latin Modern Roman}
\usepackage{amsmath,amssymb,mathtools}
\usepackage{unicode-math}
\setmathfont{Latin Modern Math}
\usepackage{xurl,hyperref}
\hypersetup{colorlinks=true,urlcolor=blue}
\setcounter{secnumdepth}{0}
\setlength{\parindent}{0pt}
\setlength{\parskip}{5pt}
\setlength{\emergencystretch}{3em}
\begin{document}
\section*{\texorpdfstring{Hilbert's Third Problem for Spherical and Hyperbolic Scissors Congruence}{Hilbert's Third Problem for Spherical and Hyperbolic Scissors Congruence}}\label{369}
\subsection{Definitions and mathematical statement}
Two three-dimensional polyhedra in $X=\mathbb H^3$ or $S^3$ are scissors congruent if finite decompositions into polyhedral pieces can be matched by isometries. Their classes form the scissors group $\mathcal P(X)$, imposing additivity under cuts. The Dehn invariant is the additive invariant represented, in the source's angular convention, by
\[
 D(P)=\sum_{e}\ell(e)\otimes(\theta_e\bmod\pi{\mathbb{Q}}).
\]
Here $\ell(e)$ and $\theta_e$ are edge lengths and dihedral angles; the precise target group and spherical relations are those of Neumann's formulation. The task, separately for both geometries, is injectivity of volume on $\ker D$. Equivalently, equal volume and equal Dehn invariant should imply scissors congruence. In the spherical scissors group volume uses its stated full-sphere convention; one must not replace the group or allow ideal pieces without checking those conventions.
\par\smallskip\textit{Formulation and concept references:} \href{https://arxiv.org/pdf/math/9712226}{[S1]}.

\subsection{Short English statement}
In spherical and hyperbolic three-dimensional geometry, do volume and the Dehn invariant together completely determine scissors congruence?
\subsection{Research attempt: exact spherical grids and a compact hyperbolic tile}
\textbf{Status.} Injectivity of volume on the full Dehn kernel, in either geometry, is \textbf{UNRESOLVED} by this attempt. We prove the necessity of the invariants directly, construct spherical families with an explicit cutting algorithm, and construct a compact right-angled hyperbolic polyhedron giving a nonzero Dehn-kernel class. The essential remaining issue is relations among general kernel classes, not additivity of the displayed formula.

\textbf{Source and convention audit.} Neumann [S1, Conjecture 1.3] states the two injectivity assertions. His Section 1 uses actual real volume; later regulator descriptions quotient the spherical group by sphere or rational-suspension classes. Those quotients cannot replace the original spherical group without tracking the lost volume. A more recent primary paper of Campbell--Zakharevich [E1] reduces generalized scissors questions to injectivity of Cheeger--Chern--Simons invariants; its abstract describes a reduction, not a general solution of the regulator-injectivity problem. We use only compact pieces below, and do not move to ideal hyperbolic tetrahedra without proving a convention-preserving bridge.

\subsubsection{1. Additivity of the Dehn invariant under a cut}
Consider first cutting a polyhedron by one totally geodesic plane. Existing edges which are subdivided contribute the same total length with the same angle. If a cut plane contains an original edge, its old dihedral angle is split into the sum of the new angles, so bilinearity preserves its contribution. A newly created edge in the interior of an old face has two complementary dihedral angles whose sum is $\pi$; its total contribution is
\[
 \ell\otimes(\theta+(\pi-\theta)\bmod\pi\mathbb Q)=0.
\]
If several pieces meet at a new internal edge, the angles sum to $2\pi$, again zero in the angular quotient. Refining any finite polyhedral decomposition by its face planes reduces additivity to these local calculations. Therefore $D$ is unchanged by finite cuts and reassembly by isometries. Volume is likewise additive because common faces have three-dimensional measure zero.

Thus equal volume and Dehn invariant are necessary for scissors congruence. The converse asks whether every formal difference annihilated by both invariants is generated by cutting relations. Defining the invariants does not prove that kernel statement.

The quotient is by $\pi\mathbb Q$, not just by $\pi\mathbb Z$. In particular every rational multiple of $\pi$ has zero angular class. Any polyhedron all of whose dihedral angles are rational multiples of $\pi$ has $D=0$, regardless of whether its edge lengths are rational or algebraic. This supplies many candidates in the kernel but no classification of their classes.

\subsubsection{2. A concrete spherical tetrahedron from two circular arcs}
Identify the unit three-sphere with
$\{(z_1,z_2)\in\mathbb C^2:|z_1|^2+|z_2|^2=1\}$. Write
\[
 z_1=\cos\eta\,e^{i\alpha},\qquad
 z_2=\sin\eta\,e^{i\beta},\qquad0\le\eta\le\pi/2.
\]
For intervals $I=[\alpha_0,\alpha_0+A]$ and
$J=[\beta_0,\beta_0+B]$ with $0<A,B<\pi$, let $P(I,J)$ be the region with $\alpha\in I$, $\beta\in J$ and the full indicated range of $\eta$.

This is the spherical join of the two arcs in the orthogonal great circles $z_2=0$ and $z_1=0$. It is a genuine spherical tetrahedron with the four arc endpoints as vertices. Its bounding surfaces $\alpha=\text{constant}$ and $\beta=\text{constant}$ lie in real linear hyperplanes of $\mathbb R^4$, so its faces are totally geodesic. The restrictions $A,B<\pi$ make the sectors convex and avoid ambiguous long-arc pieces.

The induced metric is
\[
 d\eta^2+\cos^2\eta\,d\alpha^2+\sin^2\eta\,d\beta^2,
\]
so direct integration gives
\[
 \operatorname{vol}P(I,J)
 =\int_0^{\pi/2}\sin\eta\cos\eta\,d\eta\int_I d\alpha\int_Jd\beta
 =\frac{AB}{2}.                                      \tag{1}
\]
Integrating over both full circles yields $\operatorname{vol}(S^3)=2\pi^2$, fixing the full-sphere convention.

The edge in the first circle has length $A$ and dihedral angle $B$: in its normal two-plane the region is precisely the $\beta$-sector. The second-circle edge has length $B$ and angle $A$. Each of the four joining edges has length $\pi/2$, since the circles lie in orthogonal two-planes, and has dihedral angle $\pi/2$. Therefore
\[
 D(P(I,J))=A\otimes(B\bmod\pi\mathbb Q)
          +B\otimes(A\bmod\pi\mathbb Q).              \tag{2}
\]
In particular rational angular widths $A/\pi,B/\pi\in\mathbb Q$ give $D=0$. Formula (2) does not declare all such joins trivial: their volumes in (1) are positive.

\subsubsection{3. A finite cutting algorithm for all rational rectangular grids}
Partition the $\alpha$-circle into $m\ge3$ equal arcs and the $\beta$-circle into $n\ge3$ equal arcs. Their joins give $mn$ tetrahedra with disjoint interiors covering $S^3$. Every cell is congruent, since independent rotations of $z_1$ and $z_2$ move any pair of arcs to any other pair. Each has
\[
 A=2\pi/m,\quad B=2\pi/n,\quad D=0,\quad
 \operatorname{vol}=\frac{2\pi^2}{mn}.                 \tag{3}
\]
Let $P,Q$ be finite unions of cells, allowing finite disjoint unions as in scissors bookkeeping. If they use the same grid and have equal volume, they have the same number of cells. Match cells bijectively and use the explicit rotations to reassemble one union into the other. This is an actual finite dissection, not merely a group equation.

The algorithm extends to two different rational grids and, more generally, to finite unions of joins whose arc endpoints have rational multiples of $2\pi$ as coordinates. Choose a common denominator $N$ for all $\alpha$-endpoints and a common denominator $M$ for all $\beta$-endpoints, increasing both to at least three. Cut along those grid planes. Both polyhedra become unions of congruent $2\pi/N$ by $2\pi/M$ join cells. Equality of volume then gives equality of their integer cell counts and hence a matching.

In the scissors group, these classes form a rational line generated by the full sphere. The relation $mn[P_{m,n}]=[S^3]$ follows from the explicit partition. Any two rational grid descriptions can be compared after a common refinement, so their classes are determined by the rational multiple of $2\pi^2$ which is their volume. Thus volume is injective on this concrete subspace of $\ker D$.

For example, a $4$ by $4$ grid has sixteen orthant tetrahedra, each with volume $\pi^2/8$. Refining to a $12$ by $20$ grid divides each into fifteen cells; the total counts and volumes agree exactly. The saved diagnostic checks common-denominator refinements over many small choices of grid dimensions using rational arithmetic.

This does not show that every spherical polyhedron with rational dihedral angles admits such a grid dissection. Its vertices and faces need not lie in two orthogonal circular families. That unproved extension would be a substantial scissors-congruence theorem, not a harmless coordinate choice.

\subsubsection{4. An explicit compact right-angled hyperbolic dodecahedron}
We next construct a compact polyhedron in $\mathbb H^3$ with $D=0$, keeping all vertices away from the ideal boundary. Use the Klein unit-ball model and a Euclidean regular dodecahedron centered at zero with circumradius one. Its vertices can be obtained by normalizing by $\sqrt3$ the twenty vectors
\[
 (\pm1,\pm1,\pm1),\quad
 (0,\pm\varphi^{-1},\pm\varphi),\quad
 (\pm\varphi^{-1},\pm\varphi,0),\quad
 (\pm\varphi,0,\pm\varphi^{-1}),
\]
where $\varphi=(1+\sqrt5)/2$. Each has squared norm three before normalization.

Write its face inequalities as $n_i\cdot x\le h$, with outward Euclidean unit normals $n_i$. A direct computation from these vertices and the dual icosahedral normals gives
\[
 h^2=\frac{5+2\sqrt5}{15},\qquad
 n_i\cdot n_j=c=\frac1{\sqrt5}
 \quad\text{for adjacent faces}.                     \tag{4}
\]
For instance a face normal proportional to $(0,\varphi,1)$ has support value
$(\varphi+1)/(\sqrt3\sqrt{\varphi^2+1})$, whose square is the displayed $h^2$; neighboring such normals have normalized dot product $c$.

Scale the vertices by $r\in(0,1)$ inside the Klein ball. The face planes become $n_i\cdot x=rh$. In the hyperboloid model with Lorentz form $-t^2+|x|^2$, their outward spacelike normal vectors are $(rh,n_i)$, of squared norm $1-r^2h^2$. The hyperbolic interior dihedral angle $\theta$ therefore satisfies
\[
 \cos\theta=-\frac{\langle(rh,n_i),(rh,n_j)\rangle_{\mathrm{Lor}}}
                    {1-r^2h^2}
 =\frac{r^2h^2-c}{1-r^2h^2}.                          \tag{5}
\]
Choose
\[
 r^2=\frac{c}{h^2}=3\sqrt5-6.
\]
This lies strictly between zero and one, because $2<\sqrt5<7/3$. Formula (5) then gives $\theta=\pi/2$ at every edge. All vertices have Klein norm $r<1$, so the resulting polyhedron $T$ is compact, nondegenerate and has finite strictly positive volume.

Its Dehn invariant is zero term by term because every angle is $\pi/2$. Thus $[T]\in\ker D$ is nonzero: the volume homomorphism sends it to a positive number. This construction verifies a nontrivial compact hyperbolic kernel element without ideal truncation or a numerical approximation to an angle.

If two polyhedra each admit finite decompositions into copies of this fixed $T$, equality of volume forces equal tile counts and hence an explicit scissors congruence. More generally volume is injective on the cyclic subgroup generated by $[T]$, since $a\operatorname{vol}(T)=0$ with $a\in\mathbb Z$ implies $a=0$. This is a limited but exact subcase of the desired injectivity, and does not compare different right-angled tiles with unrelated volumes.

\subsubsection{5. Tensor cancellation is algebraic, not a decimal identity}
A Dehn invariant lives in a tensor product over $\mathbb Q$; it is not obtained by multiplying real lengths and angles to a single scalar. Even if a numerical sum $\sum_e\ell(e)\theta_e$ is zero or nearly zero, that need not make the tensor sum vanish. Conversely rational-linear relations can cancel tensors in ways not visible from treating the edge terms as unrelated symbols.

The rational-angle cases above avoid that ambiguity because each angular class is exactly zero. In a more general calculation one needs exact rational relations among the lengths and among the angular classes. Given finite lists lying in specified finite-dimensional rational spaces, one can choose bases and compute a coefficient matrix for the tensor. Its vanishing is then equivalent to exact vanishing of every coefficient. But identifying those rational spaces and proving their relations for arbitrary transcendental angle and length values is not supplied by floating-point evaluation.

Similarly, arbitrarily accurate equality of computed volumes is not a certificate of exact equality. A search may suggest a relation among Dehn-kernel polytopes, but the conjecture demands an exact relation and then a finite dissection. Our grid algorithm has both: rational normalized volumes and an explicit common refinement.

\subsubsection{6. Why regulator and ideal-tetrahedron reductions need more work}
The source relates the hyperbolic kernel to a part of the Bloch group and volume to a dilogarithmic regulator. Such an identification translates the injectivity question; it does not prove the regulator injective. A finite-dimensional linear map can have a nonzero kernel, and describing its domain by generators and relations does not remove that possibility without analyzing the map.

The spherical side requires further care. The sphere class has $D=0$ but volume $2\pi^2\ne0$. Therefore it cannot be set equal to zero while retaining real volume as the same homomorphism. Quotients by sphere or rational-suspension classes lead to volume modulo the corresponding periods, as explicitly recorded later in [S1]. A proof in that quotient must recover the omitted volume information before concluding the catalogue assertion.

Ideal hyperbolic simplices introduce a different issue: their edge lengths are infinite, and the compact formula must be replaced by a truncation-independent construction. Neumann explains how horospherical changes cancel using angular relations, but a compact polyhedron cannot simply be declared a sum of ideal simplices within the original scissors group without the relevant theorem and conventions. Our right-angled tile avoids this passage entirely.

Finally, the Euclidean sufficiency theorem cannot be transported by rescaling a curved polyhedron. Spherical and hyperbolic spaces have fixed curvature and lack the Euclidean similarity transformations which scale every length while preserving angles. Infinitesimal agreement with Euclidean geometry yields approximations, not finite exact cutting identities.

\subsubsection{7. Verification and the unresolved kernel}
The exact checks verify the rational-grid volumes and refinements, the golden-ratio identities in (4), and $r^2h^2=c$ in (5). The hyperbolic polyhedron remains compact because $r<1$ is checked symbolically. These computations validate the constructed pieces and the restricted matching algorithm; they do not enumerate the whole scissors group.

The proved subcases are the rational spherical join-grid family and equal-copy assemblies of a fixed compact hyperbolic tile. The first remaining step is to show that an arbitrary class with both zero Dehn invariant and zero real volume is generated by actual cutting relations, separately in each geometry. None of the explicit families spans that unknown kernel, and a regulator reformulation alone does not show injectivity. Both general assertions remain \textbf{UNRESOLVED} by this attempt.

\subsection{Sources}
\begin{itemize}
\item[S1]Neumann, Hilbert's 3rd Problem and invariants of 3-manifolds, Geom. Topol. Monogr. 1 (1998), 383-411.
\url{https://arxiv.org/pdf/math/9712226}.
\textit{Formulation source.}
\item[E1]J. A. Campbell and I. Zakharevich, \emph{Hilbert's third problem and a conjecture of Goncharov}, Advances in Mathematics 451 (2024), 109757.
\url{https://doi.org/10.1016/j.aim.2024.109757}.
\textit{Primary abstract and introduction inspected 27 September 2026; the reduction to regulator injectivity is distinguished from a proof of that injectivity.}
\end{itemize}
\end{document}
